%==============================================================================
% APPENDIX A — BENCHMARK DATASET
%==============================================================================
\chapter{Benchmark Dataset: Question-Answer Pairs}
\label{app:benchmark}

This appendix documents the structure and annotation guidelines for the 50-item benchmark dataset used in the evaluation of Chapter~\ref{chap:evaluation}. The full dataset is provided in machine-readable format (JSON Lines) alongside this thesis.

\section*{Dataset Structure}

Each record in the dataset has the following structure:

\begin{lstlisting}[language=Python, caption={Benchmark dataset record schema}]
{
  "id": "Q001",
  "category": "factoid" | "multi-hop" | "procedural" | "unanswerable",
  "query": "Natural language question text",
  "ground_truth_answer": "Expected answer text",
  "relevant_chunk_ids": ["chunk_id_1", "chunk_id_2"],
  "source_document_titles": ["Document Title 1"],
  "annotator_1_label": 1,
  "annotator_2_label": 1,
  "cohen_kappa_pair": 0.85,
  "notes": "Optional annotation notes"
}
\end{lstlisting}

\section*{Annotation Guidelines}

Annotators were provided the following guidelines:

\begin{enumerate}
  \item Read the question and the candidate document chunk.
  \item Label the chunk as \textbf{1 (Relevant)} if it contains information necessary or sufficient to answer the question, or \textbf{0 (Not Relevant)} if it does not.
  \item A chunk may be labeled Relevant even if it provides only partial information (relevant to multi-hop synthesis).
  \item For unanswerable questions: all chunks should be labeled 0 if the question cannot be answered from the corpus.
  \item When in doubt, label 0 and note the ambiguity.
\end{enumerate}

\textit{Note: The DNSI-specific document content and individual Q\&A pairs are omitted from this appendix due to the confidentiality agreement. The anonymized synthetic test corpus is provided in the digital supplementary materials.}
